\begin{center}
{\Huge \scshape Qiong Liu, Ph.D.} \\ \vspace{4pt}
\small
\faPhone\ 475-280-1364 \quad
\href{mailto:qiongliu.work@gmail.com}{qiongliu.work@gmail.com} \quad
\href{https://www.linkedin.com/in/qiong-l-529b23196}{LinkedIn} \quad
\href{https://scholar.google.com/citations?user=SvnpHCkAAAAJ&hl=en}{Google Scholar}

Authorized to work in the U.S. without sponsorship
\end{center}

\section{Professional Summary}
\textbf{Machine Learning / Algorithm Engineer} with experience developing data-driven models for \textbf{noisy, high-dimensional, and temporally structured data}. Strong background in \textbf{signal processing, self-supervised learning, anomaly detection, and spatiotemporal modeling} for real-world systems. Experienced in building robust pipelines for \textbf{feature extraction, model training, and quantitative evaluation} under limited and imperfect data conditions. Skilled at translating complex data patterns into actionable insights for system optimization.

\section{Core Expertise}
\textbf{Machine Learning:} Representation learning, self-supervised learning, anomaly detection, contrastive learning \\
\textbf{Signal \& Data Modeling:} Time-series analysis, spatiotemporal modeling, signal decomposition, noise modeling \\
\textbf{Computer Vision:} Image restoration, defect detection, pattern recognition, feature extraction \\
\textbf{Programming:} Python, PyTorch, NumPy, SciPy, MATLAB \\
\textbf{Engineering Strengths:} Model robustness, generalization, data-driven system optimization

\section{Experience}

\textbf{Canon Medical Research USA} \hfill Apr 2025 -- Present \\
\textit{Reconstruction Scientist, Vernon Hills, IL}
\begin{itemize}[leftmargin=*]
\item Developed methods to \textbf{extract dual motion signals} from noisy dynamic data with overlapping temporal patterns and distribution shift
\item Built \textbf{signal-aware motion correction models} to address noise--feature trade-offs using domain-informed constraints
\item Developed \textbf{diffusion-based denoising models} and integrated them into reconstruction pipelines to improve overall data quality
\end{itemize}

\textbf{Canon Medical Research USA} \hfill Jun 2023 -- May 2024 \\
\textit{Research Scientist Intern, Vernon Hills, IL}
\begin{itemize}[leftmargin=*]
\item Developed \textbf{efficient attention-based denoising models} with modified transformer architectures for faster inference
\item Designed \textbf{physics-informed noise modeling methods} for signal-dependent noise in low-SNR data
\item Built deep learning models for \textbf{scatter estimation}, reducing a major runtime bottleneck in the reconstruction pipeline
\end{itemize}

\textbf{Yale University, Yale PET Center} \hfill Aug 2020 -- May 2025 \\
\textit{Ph.D. Researcher (Machine Learning / Imaging Systems), New Haven, CT}
\begin{itemize}[leftmargin=*]
\item Developed learning-based methods for \textbf{low-SNR signal reconstruction and denoising}, using self-supervised learning and domain-informed constraints to preserve meaningful structure
\item Built robust learning pipelines for \textbf{data-scarce and heterogeneous settings}, including federated learning, fine-tuning, and adaptation across varying data distributions
\item Developed \textbf{dynamic system modeling} methods for temporal signal decomposition, quantitative analysis, and data-driven modeling of novel tracer behavior
\end{itemize}

\textbf{United Imaging Healthcare} \hfill May 2019 -- Aug 2019 \\
\textit{Student Intern, Wuhan, China}
\begin{itemize}[leftmargin=*]
\item Supported development and testing of a \textbf{surgical robotics system}, including image registration, system validation, and performance evaluation
\end{itemize}

\section{Patents}
\textbf{Edge-Guided Deformation Field Fusion for Motion Correction} \hfill 2026 \\
Patent Pending, \textit{Lead Inventor}
\begin{itemize}[leftmargin=*]
\item Proposed a spatially adaptive deformable modeling framework to improve alignment accuracy and robustness in challenging image regions
\end{itemize}

\textbf{Automatic Data-Driven Multi-Signal Decomposition} \hfill 2025 \\
Patent Pending, \textit{Lead Inventor}
\begin{itemize}[leftmargin=*]
\item Developed an AI-based framework to separate overlapping temporal signals from dynamic image sequences
\end{itemize}

\section{Selected Projects}

\textbf{Self-Supervised Denoising with Physics-Guided Constraints}
\begin{itemize}[leftmargin=*]
\item Developed a self-supervised framework for \textbf{low-SNR signal denoising} using multi-frame consistency and reconstruction objectives
\item Introduced \textbf{physics-guided constraints} (e.g., kinetic consistency) to preserve meaningful signal behavior during denoising
\item Demonstrated improved robustness and quantitative consistency compared to supervised and unguided methods
\end{itemize}

\textbf{Data-Driven Prior Modeling for Low-Signal Reconstruction}
\begin{itemize}[leftmargin=*]
\item Developed population-based and data-driven priors for \textbf{signal reconstruction under extreme noise and limited data}
\item Analyzed how \textbf{training data distribution and prior design} impact model behavior and reconstruction quality
\item Improved reconstruction fidelity while maintaining downstream quantitative accuracy
\end{itemize}

\textbf{Dynamic System Modeling for Quantitative Signal Analysis}
\begin{itemize}[leftmargin=*]
\item Developed \textbf{parametric modeling and system identification methods} for dynamic processes in sequential data
\item Applied modeling techniques to extract \textbf{quantitative features} for analysis and evaluation of complex systems
\item Enabled improved interpretability and decision-making through data-driven modeling of dynamic behavior
\end{itemize}

\section{Education}
\textbf{Yale University} \hfill Aug 2020 -- May 2025 \\
Ph.D. in Biomedical Engineering \\
Thesis: Improving Image Quality and Quantification Accuracy for Static and Dynamic PET

\textbf{Huazhong University of Science and Technology} \hfill Sep 2016 -- Jun 2020 \\
B.S. in Biomedical Engineering, Minor in Computer Science \\
GPA: 3.94/4.00

\section{Selected Publications}

\textbf{First Author}
\begin{itemize}[leftmargin=*]
\item Liu Q., et al., \textit{Parametric cardiac imaging with $^{18}$F-flutemetamol PET to evaluate the impact of tafamidis in patients with transthyretin cardiac amyloidosis}, \textbf{Journal of Nuclear Medicine}, 2026
\item Liu Q., et al., \textit{Patlak-guided self-supervised learning for dynamic PET denoising}, \textbf{IEEE Transactions on Radiation and Plasma Medical Sciences}, 2026
\item Liu Q., et al., \textit{Population-based deep image prior for dynamic PET denoising: a data-driven approach to improve parametric quantification}, \textbf{Medical Image Analysis}, 2024
\item Liu Q., et al., \textit{A personalized deep learning denoising strategy for low-count PET images}, \textbf{Physics in Medicine \& Biology}, 2022
\end{itemize}

\textbf{Co-Author}
\begin{itemize}[leftmargin=*]
\item Xia M., Xie H., Liu Q., et al., \textit{LeqMod: adaptable lesion-quantification-consistent modulation for deep learning low-count PET image denoising}, IEEE Transactions on Medical Imaging, 2025
\item Zhou B., Xie H., Liu Q., et al., \textit{FedFTN: personalized federated learning with deep feature transformation network for multi-institutional low-count PET denoising}, Medical Image Analysis, 2023
\end{itemize}

\section{Awards}
\textbf{Society of Nuclear Medicine and Molecular Imaging (SNMMI)}
\begin{itemize}[leftmargin=*]
\item Young Investigator Award (1st Place), 2025
\item Poster Award (2nd Place), 2024
\item Young Investigator Award (2nd Place), 2023
\end{itemize}

\textbf{Mitacs}
\begin{itemize}[leftmargin=*]
\item Globalink Research Award, 2019
\end{itemize}